\begin{table}[t]
\caption{Frozen controller $\times$ fault comparison, seeds 12--17. Six clusters share both bearings and all conditions. L/G: paired healthy-to-fault losses/gains out of 12; healthy survivors at risk are four/six under original/gated laws. Healthy rows reuse preserved records.}
\label{tab:steering}
\centering\small
\setlength{\tabcolsep}{3pt}
\begin{tabular}{llrrrr}
\toprule
Law & Network & $-30^\circ$ & $+30^\circ$ & Total & L/G\\
\midrule
Original & Healthy & 4/6 & 0/6 & 4/12 & ---\\
Original & 25\% deletion & 4/6 & 0/6 & 4/12 & 1/1\\
Original & 50\% deletion & 0/6 & 0/6 & 0/12 & 4/0\\
Pre-cue gate & Healthy & 6/6 & 0/6 & 6/12 & ---\\
Pre-cue gate & 25\% deletion & 3/6 & 0/6 & 3/12 & 3/0\\
Pre-cue gate & 50\% deletion & 0/6 & 0/6 & 0/12 & 6/0\\
\bottomrule\end{tabular}\end{table}

The 72-record matrix directly pairs each fault with its
same-law healthy reference (Table~\ref{tab:steering}). The healthy gate advantage
is two episodes out of 12 (+16.7 percentage points); its six-cluster conditional
range includes zero. Offset-only and no-steering healthy controls triggered
bursts in 11/12 and 12/12 cases but yielded 0/12 survival (Supplementary
Tables~S12--S14). These compare complete feedback laws, which can change sensory
trajectories and trigger timing.

At 25\% deletion, gated-minus-original net task loss was +25.0 percentage points (conditional 95\% range 8.3 to 41.7). At 50\% deletion, gated-minus-original net task loss was +16.7 percentage points (conditional 95\% range 0.0 to 33.3). The 25\% range excludes zero, but only 3 seed clusters have nonzero interactions; a post-hoc symmetric-null sign-flip sensitivity gives $p=0.25$ (Supplementary Section~13). At 50\% deletion, both controllers yield 0/12; their net-loss contrast reflects different healthy baselines. Loss and gain counts are retained separately; a healthy success difference alone is not a fault-tolerance comparison.

Among 48 faulty episodes, 41 collided; none of these collisions was flagged at any tested $c=1.5,2,3,4$ over 200--500\,ms. At 500\,ms, $c=1.5$ instead flagged 2 surviving fault episodes. This activity-only criterion did not distinguish the observed collisions under these settings; it does not establish that spike count is generally uninformative. Supplementary Table~S22 and the threshold CSV report the full sensitivity.
